% ===== PRÉAMBULE CANONIQUE — à coller VERBATIM en tête de CHAQUE .tex =====
\documentclass[11pt,a4paper]{article}
\usepackage[utf8]{inputenc}\usepackage[T1]{fontenc}\usepackage{lmodern}
\usepackage[french]{babel}
\usepackage{amsmath,amssymb,mathtools}\usepackage{icomma}
\usepackage[version=4]{mhchem}          % \ce{...} : équations & formules
\usepackage{siunitx}\sisetup{locale=FR} % \qty{}{}, \si{}, \ang{}
\usepackage{chemfig}                    % \chemfig{...} : structures développées
\usepackage{xcolor}\usepackage{colortbl}
\usepackage{tikz}\usepackage{pgfplots}\pgfplotsset{compat=1.18}
\usepackage[most]{tcolorbox}\usepackage{enumitem}\usepackage{array,booktabs}
\usepackage[a4paper,margin=2.2cm]{geometry}
\usetikzlibrary{arrows.meta,calc,angles,quotes,positioning,patterns,backgrounds,shapes.geometric}
\definecolor{primary}{RGB}{222,49,99}\definecolor{primarydark}{RGB}{150,22,55}
\definecolor{defcol}{RGB}{0,114,178}\definecolor{methcol}{RGB}{52,92,148}
\definecolor{excol}{RGB}{0,135,90}\definecolor{attncol}{RGB}{200,40,55}
\tcbset{boxbase/.style={enhanced,breakable,boxrule=0.9pt,arc=2.5pt,
  left=3mm,right=3mm,top=2.3mm,bottom=2.3mm,fonttitle=\bfseries,coltitle=white}}
% --- boîtes TUTORIEL (noms EXACTS, ne pas renommer) ---
\newtcolorbox{cle}[1][]{boxbase,colback=defcol!6,colframe=defcol,
  title={\textsc{À retenir}\ifx\relax#1\relax\relax\else\textmd{ : \itshape #1}\fi}}
\newtcolorbox{methode}[1][]{boxbase,colback=methcol!7,colframe=methcol,sharp corners,
  title={\textsc{Méthode}\ifx\relax#1\relax\relax\else\textmd{ --- \itshape #1}\fi}}
\newtcolorbox{exemplebox}[1][]{boxbase,colback=excol!5,colframe=excol,
  title={\textsc{Exemple}\ifx\relax#1\relax\relax\else\textmd{ : \itshape #1}\fi}}
\newtcolorbox{piege}[1][]{boxbase,colback=attncol!6,colframe=attncol,
  title={\textsc{Piège}\ifx\relax#1\relax\relax\else\textmd{ : \itshape #1}\fi}}
\newtcolorbox{remarque}[1][]{boxbase,colback=black!4,colframe=black!45,
  title={\textsc{Remarque}\ifx\relax#1\relax\relax\else\textmd{ : \itshape #1}\fi}}
% --- boîtes EXERCICE / CORRIGÉ (noms EXACTS, auto-comptées) ---
\newtcolorbox[auto counter]{exo}[1][]{boxbase,colback=primary!4,colframe=primary,
  title={\textsc{Exercice}~\thetcbcounter\ifx\relax#1\relax\relax\else\textmd{ --- \itshape #1}\fi}}
\newtcolorbox[auto counter]{corrige}[1][]{boxbase,colback=excol!4,colframe=excol,
  title={\textsc{Corrigé}~\thetcbcounter\ifx\relax#1\relax\relax\else\textmd{ --- \itshape #1}\fi}}
\newcommand{\R}{\mathbb{R}}\newcommand{\N}{\mathbb{N}}\newcommand{\Z}{\mathbb{Z}}\newcommand{\ds}{\displaystyle}
% =========================================================================
\begin{document}
\sisetup{per-mode=symbol}
\begin{corrige}[{convertir entre $[\ce{OH-}]$ et pH}]
\begin{enumerate}[label=\alph*)]
  \item $\mathrm{pH}=14+\log(\num{1.0e-3})=14-3=11$.
  \item $\mathrm{pH}=14+\log(\num{1.0e-6})=14-6=8$.
  \item $[\ce{OH-}]=10^{10-14}=10^{-4}=\qty{1.0e-4}{\mole\per\litre}$.
  \item $[\ce{OH-}]=10^{12-14}=10^{-2}=\qty{1.0e-2}{\mole\per\litre}$.
\end{enumerate}
\end{corrige}

\begin{corrige}[classer les sels dans les trois cas]
\begin{enumerate}[label=\alph*)]
  \item \ce{Mg(OH)2} : \textbf{cas 1} (hydroxyde) — précipite en milieu basique, se dissout en acide.
  \item \ce{CaCO3} : \textbf{cas 2} (sel d'acide faible, \ce{CO3^{2-}}) — plus soluble en milieu acide.
  \item \ce{AgCl} : \textbf{cas 3} (anion d'acide fort) — pH sans effet.
  \item \ce{FeS} : \textbf{cas 2} (sel d'acide faible, \ce{S^{2-}}) — plus soluble en milieu acide.
  \item \ce{BaSO4} : \textbf{cas 3} (anion d'acide fort) — pH sans effet.
\end{enumerate}
\end{corrige}

\begin{corrige}[vrai ou faux : sens de l'effet du pH]
\begin{enumerate}[label=\alph*)]
  \item \textbf{Vrai.} pH $\uparrow \Rightarrow [\ce{OH-}]\uparrow \Rightarrow$ l'hydroxyde précipite.
  \item \textbf{Faux.} \ce{CaCO3} (sel d'acide faible) est plus soluble en milieu \emph{acide}.
  \item \textbf{Faux.} \ce{AgCl} est du cas 3 : sa solubilité est \emph{indépendante} du pH.
  \item \textbf{Vrai.} En milieu acide, $[\ce{OH-}]$ est minuscule, l'hydroxyde se dissout davantage.
\end{enumerate}
\end{corrige}

\begin{corrige}[le pH décide-t-il la précipitation ?]
On calcule $[\ce{OH-}]=10^{\mathrm{pH}-14}$ puis $Q$, comparé à $K_{\mathrm{ps}}=\num{9.0e-12}$.
\begin{enumerate}[label=\alph*)]
  \item $[\ce{OH-}]=10^{-4}$, $Q=(0{,}010)(\num{1.0e-4})^{2}=\num{1.0e-10} > K_{\mathrm{ps}}$ :
        \textbf{précipitation} de \ce{Mg(OH)2}.
  \item $[\ce{OH-}]=10^{-6}$, $Q=(0{,}010)(\num{1.0e-6})^{2}=\num{1.0e-14} < K_{\mathrm{ps}}$ :
        \textbf{pas} de précipité.
\end{enumerate}
\end{corrige}

\begin{corrige}[quels sels sont sensibles au pH ?]
Sont \textbf{sensibles} au pH :
\begin{itemize}[leftmargin=1.6em,topsep=2pt,itemsep=1pt]
  \item \ce{Fe(OH)3} et \ce{Al(OH)3} : ce sont des \textbf{hydroxydes} (cas 1) ;
  \item \ce{CaCO3} : \textbf{sel d'acide faible} (cas 2, \ce{CO3^{2-}}).
\end{itemize}
Sont \textbf{insensibles} : \ce{AgCl} et \ce{BaSO4} (cas 3, anions d'acides forts).
\end{corrige}
\end{document}
